\documentclass[10pt]{article}

\usepackage[utf8]{inputenc}

\usepackage[T1]{fontenc}

\usepackage{amsmath}

\usepackage{amsfonts}

\usepackage{amssymb}

\usepackage[version=4]{mhchem}

\usepackage{stmaryrd}



\begin{document}

Лит.: [1] Б л и с с Г. А., Лекции по вариационному исчислению, пер. с англ., М., 1950; [2] Л а в р е н т в е в M. А., изт. М. - Л. 1950 ., К. Б. Вапнярский.



РЕГУЛЯР̈НО КОЛЬЦО в к о м м у т а т И в н о Й а л г е б р е - нётерово кольцо $A$, все локализации $A_{\downarrow}$ к-рого регулярны; здесь $\mathfrak{p}-$ простой идеал в $A$. При этом јокальное нётерово кольцо $A$ с максимальным идеалом $\mathfrak{m}$ наз. р е г у л я р ны м, если $\mathfrak{m}$ порождается $n$ әлементами, где $n=\operatorname{dim} A$, т. е. если касательное пространство $\mathfrak{m} / \mathfrak{n t}^{2}$ (как векторное пространство над полем вычетов) имеет размерность, равную $\operatorname{dim} A$. Это равносильно отсутствию особенностей у схемы Spec $A$. Локальное Р. к. $A$ всегда целостно и нормально, а также факториально (т е о р е м а А у с л е н д е р а Б у к с б а у м а), глубина его равна $\operatorname{dim} A$. Асcoциированное градуированное кольцо



\[

G_{\mathfrak{m}}(A)=\bigoplus_{i \geqslant 0} \mathfrak{m}^{i} / \mathfrak{m}^{i+1}

\]



изоморфно кольцу многочленов $k\left[X_{1}, \ldots, X_{n}\right]$. Локальное нётерово кольцо $A$ регулярно тогда и только тогда, когда регулярно его пополнение $\widehat{A} ;$ вообще, если $A \subset B$ - плоское расширение локальных колец и $B$ регулярно, то и $A$ регулярно. Для полных локальных Р. к. имеет место с т р у к т у р н а я т е о р е м а К о ә н а: они имеют вид $R\left[\left[X_{1}, \ldots, X_{n}\right]\right]$, где $R-$ поле или кольцо дискретного нормирования. Јюбой модуль конөчного типа над локальным Р. к. обладает конечной свободной резольвентой (см. Гильберта теорема о сизигиях); верно и обратное (см. [2]).



Р. к. являются любое поле и любое дедекиндово кольцо. Если $A$ регулярно, то регулярно кольцо многочленов $A\left[X_{1}, \ldots, X_{n}\right]$ и кольцо формальных степенных рядов $A\left[\left[X_{1}, \ldots, X_{n}\right]\right]$ над $A$. Если $a \in A-$ необратимый элемент локального $\mathrm{P}$. к., то $A / a A$ регулярно тогда и только тогда, когда $a \notin \mathfrak{m}^{2}$.



Jum.: [1] 3 а р и с с к и й О., С а м ю э ль П., Коммутативная алгебра, пер. с англ., т. 2, М., 1963; [2] С е p p Ж.- П., "Математика", 1963 , т. 7, № 5, c. 3-93; [3] G r o th e n d ie c k A., D i e u d o n n é J. (red.), Eléments de géométrie algébrique, chap. 4, Р., 1964. В. И. Данилов.



РЕГУЛЯРНОЕ КОЛЬЦО (в смысле Неймана) ассоциативное кольцо (обычно с единицей), в к-ром уравнение $a x a=a$ разрешимо для любого $a$. Следующие своїства ассоциативного кольца $R$ с единицей равносильны: a) $R$ есть Р . к.; б) каждый главный левый идеал кольца $R$ порождается идемпотентом; в) главные левые идеалы кольца $R$ образуют подрешетку в решетке всех его левых идеалов, являющуюся дедекиндовой решеткой с дополнениями; г) каждый главный левый идеал кольца $R$ имеет дополнение в структуре всех его левых идеалов; д) все левые $R$-модули плоские; е) имеют место правые аналоги свойств б) - д) (см. [3], [4], $[5],[8],[10])$. Ввиду д) Р. к. иногда наз. а б с о л ю т н о п л о к и м и. Коммутативное кольцо регулярно тогда и только тогда, когда инъективны все простые модули над ним (см. [5]). Любой конечно порожденный левый (правый) идеал Р. к. оказывается главным и, следовательно, выделяется прямым слагаемым. Всякий неделитель нуля в Р.к. обратим. Радикал Джекобсона Р . к. равен нулю. Кольцо матриц над Р . к. оказывается Р. к. Класс Р. к. замкнут относительно перехода к прямым произведениям и факторкольцам. Идеал Р. к. является Р. к. (возможно, без единицы). Если Р. к. нётерово или совершенно (слева или справа), то оно оказывается классически полупростым кольцом. Всякое классически полупростое кольцо регулярно. Более того, регуля рным оказывается кольцо эндоморфизмов векторного пространства над телом (даже бесконечномерного), а также факторкольцо кольца эндоморфизмов инъективного левого (правого) модуля над любым кольцом по его радикалу Джекобсона (см. [3]) . В частности, всякое самоиньективное слева (справа) кольцо с нулевым радикалом Дґекобсона регулярно. Групповое кольцо групшы $G$ над Р. к. регулярно тогда и только тогда, когда любая конечно порожденная подгруппа в $G$ конечна и порядок каждой из таких подгрупп обратим в исходном Р. к. (см. [3]). Кольца эндоморфизмов всех свободных левых $R$-модулей регулярны в том и только в том случае, когда $R$ классически полупросто [6]. Счетно порожденные односторонние идеалы Р.к. проективны [8].



Если $R$ есть Р. к., то конечно порожденные подмодули левого $R$-модуля $R^{n} n$-мерных строк над $R$ образуют дедекиндову решетку $L$ с дополнениями, являющуюся подрешеткой решгтки всех подмодулей модуля $R^{n}$. Решетка $L$ содержит 0 д н о р о д н ы й $a_{1}, \ldots, a_{n}$, то есть эти әлементы независимы (см. Дeдекиндова решетка), их сумма равна наибольшему элементу из $L$ (а именно, $\left.R^{n}\right)$ и любые $a_{i}$ и $a_{i}$ п е р с п е кт и в н ы, что означает существование для них общего дополнения. Наоборот, всякая дедекиндова решетка с дополнениями, обладающая однородным базисом, coдержащим не менее четырех элементов, изоморфна решетке $L$ для подходящего Р. к. $R$. Решетка $L$ изоморфна решетке главных левых идеалов кольца всех $(n \times n)$ матриц над $R$ (см. [4], [10]).



Важный частный случай Р.к.- с т р о г о р е г ул я р н о е к о л ц о, в к-ром, по определению, разрешимо уравнение $a^{2} x=a$. Равносильны следующие свойства Р.к. $R$ : а) $R$ строго регулярно; б) $R$ не содержит ненулевых нильпотентных элементов; в) все идөмпотенты кольца $R$ центральны; г) каждый левый (или каждый правый) идеал кольца $R$ является двусторонним; д) решетка главных левых (правых) идеалов кольца $R$ дистрибутивна; е) мультипликативная полугруппа кольца $R$ является инверсной полугруппой (см. [7], [8]).



Другой подкласс класса Р. к. образуют $u$-р е г у л я рны е ко льц а, где, по определению, уравнение $a \times a=a$ имеет в качестве решения обратимый элемент. В классе Р. к. u-регулярные кольца характеризуются транзитивностью перспективности в решетке конечно порожденных подмодулей суммы двух экземпляров основного кольца, а также возможностью сокращать прямую сумму на конечно порожденный проективный модуль (см. [4], [8]).



Р. к. наз. н е п р е р ы в ны м с л е в а, если непрерывна решетка его главных левых идеалов. Непрерывное Р. к. и-регулярно и разлагается в прямую сумму строго Р.к. и самоинъективного кольца. На Р. к. может быть определена псевдоранг-функция, являющаяся аналогом меры на булевой алгебре. Она определяет псевдометрику. Пополнение Р. К. по этой метрике оказывается самоинъективным Р. к. (см. [8]).



Р. к. являются частным случаем $л-р$ е г у л я р ны к о л е ц, в к-рых, по определению, для каждого элемента $a$ найдутся элемент $x$ и натуральное число $n$ такие, что $a^{n} x a^{n}=a^{n}$.



Как двусторонний аналог Р. к. можно рассматривать би р е гул я рны е ко л.ь ц а, в к-рых, по определению, каждый главный двусторонний идеал порождается центральным идемпотентом. Каждый двусторонний идеалг бирегулярного кольца является пересечением его максимальных двусторонних идеалов. Всякое бирегулярное кольцо с единицей изоморфоно кольцу глобальных сечений с бикомпактными носителями пучка простых колец с единицей над бикомпактным вполне несвязным хаусдорфовым топологич. пространством, и всякое такое кольцо глобальных сечений бирегулярно (см. [2]). В коммутативном случае классы бирегулярных, строго Р. к. и Р. к. совпадают, и простые кольца в последнеӥ теореме заменяются полями.



Близки к Р к. б э р о ски е ко льц а, определяемые тем условием, что каждый левый (или, что при наличии единицы равносильно, каждый правый)





\end{document}